\chapter*{Список сокращений}
\addcontentsline{toc}{chapter}{Список сокращений}
\markboth{Список сокращений}{}

\begin{longtable}{L{2.6cm}L{13.2cm}}
\textbf{АОП} & аспектно-ориентированное программирование
   (\eng{Aspect-Oriented Programming}, AOP) \\[2pt]
\textbf{БД} & база данных \\[2pt]
\textbf{ООП} & объектно-ориентированное программирование \\[2pt]
\textbf{СУБД} & система управления базами данных \\[2pt]
\textbf{API} & \eng{Application Programming Interface}~--- программный интерфейс
   приложения \\[2pt]
\textbf{CI/CD} & \eng{Continuous Integration~/ Continuous Delivery}~---
   непрерывная интеграция и непрерывная поставка \\[2pt]
\textbf{CRUD} & \eng{Create, Read, Update, Delete}~--- четыре базовые операции
   над данными \\[2pt]
\textbf{DAO} & \eng{Data Access Object}~--- объект доступа к данным \\[2pt]
\textbf{DI} & \eng{Dependency Injection}~--- внедрение зависимостей \\[2pt]
\textbf{DTO} & \eng{Data Transfer Object}~--- объект передачи данных \\[2pt]
\textbf{FXML} & язык разметки интерфейса в JavaFX \\[2pt]
\textbf{GC} & \eng{Garbage Collector}~--- сборщик мусора \\[2pt]
\textbf{HTTP} & \eng{HyperText Transfer Protocol}~--- протокол передачи
   гипертекста \\[2pt]
\textbf{IoC} & \eng{Inversion of Control}~--- инверсия управления \\[2pt]
\textbf{JAR} & \eng{Java ARchive}~--- архив с классами Java \\[2pt]
\textbf{JDBC} & \eng{Java Database Connectivity}~--- интерфейс доступа к БД \\[2pt]
\textbf{JDK} & \eng{Java Development Kit}~--- комплект разработчика \\[2pt]
\textbf{JIT} & \eng{Just-In-Time}~--- динамическая компиляция байт-кода \\[2pt]
\textbf{JNI} & \eng{Java Native Interface}~--- интерфейс вызова нативного кода \\[2pt]
\textbf{JPA} & \eng{Jakarta Persistence API}~--- спецификация ORM для Java \\[2pt]
\textbf{JRE} & \eng{Java Runtime Environment}~--- среда выполнения \\[2pt]
\textbf{JSON} & \eng{JavaScript Object Notation}~--- текстовый формат обмена
   данными \\[2pt]
\textbf{JVM} & \eng{Java Virtual Machine}~--- виртуальная машина Java \\[2pt]
\textbf{JWT} & \eng{JSON Web Token}~--- формат токена аутентификации \\[2pt]
\textbf{MVC} & \eng{Model--View--Controller}~--- модель-представление-контроллер \\[2pt]
\textbf{ORM} & \eng{Object-Relational Mapping}~--- объектно-реляционное
   отображение \\[2pt]
\textbf{POJO} & \eng{Plain Old Java Object}~--- простой объект Java без
   обязательного наследования \\[2pt]
\textbf{REST} & \eng{Representational State Transfer}~--- архитектурный стиль
   веб-сервисов \\[2pt]
\textbf{SOLID} & пять принципов проектирования классов \\[2pt]
\textbf{SQL} & \eng{Structured Query Language}~--- язык запросов к БД \\[2pt]
\textbf{UML} & \eng{Unified Modeling Language}~--- унифицированный язык
   моделирования \\[2pt]
\textbf{WAR} & \eng{Web Application Archive}~--- архив веб-приложения \\[2pt]
\textbf{XML} & \eng{eXtensible Markup Language}~--- расширяемый язык разметки \\
\end{longtable}
